\subsubsection*{Bayes' Rule Explained}



\subsubsection*{The idea}

You start with a guess about how likely something is. Then you see a clue. Bayes' rule tells you how much that clue should change your guess.

\subsubsection*{An example}

Say 1 out of every 100 students in your school has the flu (1\%). The school has a test for it, and the test is pretty good: it correctly says ``positive'' for 90\% of students who really have the flu. But it also says ``positive'' by mistake for 9\% of healthy students.

You test positive. What are the chances you actually have the flu? Most people say ``about 90\%.'' That's wrong. Let's work it out with 1,000 students:

\begin{itemize}
  \item 10 students have the flu (1\% of 1,000). The test catches 9 of them.
  \item 990 students are healthy. The test wrongly flags about 89 of them (9\% of 990).
  \item So about 98 students test positive in total (9 + 89).
  \item Of those 98, only 9 really have the flu.
\end{itemize}

Chance you have the flu, given a positive test:
\[
\frac{9}{98} \approx 9\%
\]

The test is good, but the flu is rare, so most positive results come from healthy people getting false alarms.

\subsubsection*{The formula}

Bayes' rule is just that counting written in symbols:
\[
P(A \mid B) = \frac{P(B \mid A) \times P(A)}{P(B)}
\]

\begin{itemize}
  \item $P(A \mid B)$: the chance of A, given that B happened (chance of flu, given a positive test). This is what you want.
  \item $P(B \mid A)$: the chance of B if A is true (chance of a positive test if you have the flu: 90\%).
  \item $P(A)$: how likely A was before any evidence (1\% had the flu). This is called the \emph{prior}.
  \item $P(B)$: how likely B is overall (about 9.8\% of everyone tests positive).
\end{itemize}

Plugging in:
\[
\frac{0.90 \times 0.01}{0.098} \approx 0.09 \text{, or } 9\%
\]

\subsubsection*{The takeaway}

Evidence matters, but so does what you already knew. A strong clue about a rare thing can still leave the thing unlikely. Bayes' rule is how you combine both properly.

\subsubsection*{Why 90\% Is Wrong}

The 90\% answer feels right because it mixes up two different questions:

\begin{enumerate}
  \item \textbf{If you have the flu, how likely is a positive test?} This is 90\%. It's the test's accuracy.
  \item \textbf{If you got a positive test, how likely is it that you have the flu?} This is what you actually want to know, and it's about 9\%.
\end{enumerate}

They sound almost the same, but they are not. The test's 90\% only describes what happens \emph{among sick people}. It says nothing about how many sick people there are in the first place.

\subsubsection*{Why the second question has a different answer}

Think about who ends up in the ``positive'' group. Two kinds of students are in it:

\begin{itemize}
  \item Sick students the test correctly caught: \textbf{9}
  \item Healthy students the test wrongly flagged: \textbf{about 89}
\end{itemize}

The healthy group is huge (990 students), so even a small 9\% mistake rate produces a lot of false alarms. 9\% of 990 is far more than 90\% of 10. The false alarms simply outnumber the real catches, roughly 10 to 1.

So if you pick a random student from the positive group, they are much more likely to be a healthy student with a false alarm than a sick student.

\subsubsection*{An everyday comparison}

Imagine a metal detector that beeps for 90\% of real knives, but also beeps by accident for 9\% of harmless things like belt buckles and keys. If only 1 in 100 people carries a knife, then almost every beep is a belt buckle. A beep is a clue, but not proof, because knives are rare to begin with.

\subsubsection*{Why intuition fails}

Our brains focus on the test's accuracy (the vivid, exciting number) and ignore how rare the flu is (the boring background number). This mistake is called \textbf{base rate neglect}. The ``base rate'' is how common something is before any testing, and it has a big effect on the final answer.

\textbf{Quick check:} if half the school had the flu, the same positive test would be far more convincing. The test hasn't changed, only the base rate has. That is why you can't judge a test result without knowing how common the thing is.